\scanpage{12}

et on peut demander des conditions des types suivants

(i)\deleted{\unsure{(ii)}} la restriction de $\varphi^*_C$ à $\underline{\mathrm{Hom}}_{\bar\delta}$ est fidèle (cas particulier de c)

(\unsure{$i'$|$ii$})\deleted{\unsure{(iii)}} la restriction de $\varphi^*_C$ à $\underline{\mathrm{Hom}}_{\bar\delta}$ est \deleted{\unsure{un}} surjectif sur les Hom

\inserted{(ii) $\varphi^*_C$ envoie $\underline{\mathrm{Hom}}_{\bar\delta}$ dans $\underline{\mathrm{Hom}}_\delta$ : cas particulier de c), et permet d'exprimer c)}

($ii'$)\deleted{\unsure{(iv)}} $\underline{\mathrm{Hom}}_\delta$ est dans l'image essentielle de $\underline{\mathrm{Hom}}_{\bar\delta}$

\uline{NB} (\unsure{$i'$|$i$}) \deleted{et (ii)} est dans la nature \inserted{d'une condition} d'existence de flèches (ou systèmes de flèches) ds $C$, \deleted{\unsure{et}} ($ii'$) dans la nature d'existence d'objets et de diagrammes d'objets. Ainsi, ($ii'$) généralise a) et b) [ \deleted{\unsure{(ii$'$) impli}} des conditions de type ($i'$) peuvent-elles \inserted{être} exprimées comme conditions de type ($ii'$) ? ]

\margin{? !}

Donc finalement on revient à deux types de conditions permettant d'exprimer les autres

\[
\begin{aligned}1^\circ)&\ \deleted{\unsure{\varphi^*}}\\2^\circ)&\end{aligned}
\left\{\begin{aligned}
&\mathrm{Im\,ess}(\varphi^*_C\vert\underline{\mathrm{Hom}}_{\bar\delta})\subset\underline{\mathrm{Hom}}_\delta\\
&\mathrm{Im\,ess}(\varphi^*_C\vert\underline{\mathrm{Hom}}_{\bar\delta})\supset\underline{\mathrm{Hom}}_\delta
\end{aligned}\right.
\begin{aligned}
&\text{implique conditions d'exactitude }\inserted{\text{habituelles}}\\
&\text{implique }\inserted{\text{conditions d'}}\text{existence (y compris conditionnelle) de limites}
\end{aligned}
\]

Vérifier que cel permet \inserted{bien} d'exprimer l'\unsure{injectivité|intégralité|égalité} resp. surjectivité de $\varphi^*_C\vert\underline{\mathrm{Hom}}_\delta$ sur les Hom ....

$2^\circ$) \uline{Conditions d'exactitude sur foncteurs} $\deleted{\unsure{g}}F:C\longrightarrow C'$

On a déjà vu comment on peut en exprimer, avec des paquets de \deleted{diagrammes}\inserted{cônes} de $C$. Mais on veut des conditions \deleted{\unsure{ind}} d'ex valables pour \uline{tt} $C$, sans donnée supplémentaire sur $C$, -- donc des façons d'associer à tt $C$ des paquets de \deleted{diagrammes}\inserted{cônes} de $C$. Une façon de le faire est de partir de $\Theta$, muni de $\delta\deleted{\unsure{\xi}}=(\delta',\delta'')$ et $\lambda=(\lambda',\lambda'')$ \inserted{(paquets de cônes)} et de prendre les $u:\Theta\unsure{-|\to}C$ qui sont $\delta$-exacts, et demander que $F$ soit $u(\lambda)$-exact pour tt tel $u$\deleted{\unsure{ précis}}. On peut prendre un paquet de \unsure{ce}\deleted{\unsure{$T_\alpha$}} $(\Theta_\alpha,\delta_\alpha,\lambda_\alpha)=T_\alpha$, soit $T=(T_\alpha)$

On désire que \deleted{\unsure{les}} la $T$-exactitude soit propriété

% archive sheet number: 11
